% ============================================================================
% WORKBOOK — Section 2.2: Connecting Percents and Proportions
% CCSS 7.RP.A.3
% Practice-focused reinforcement for the study guide topic
% ============================================================================
\section{Connecting Percents and Proportions}
\topicTitle{Connecting Percents and Proportions}

\begin{quickReview}{Percents as Proportions}
``Percent'' means ``per hundred.'' Every percent problem can be written as a proportion:
$$\frac{\text{part}}{\text{whole}} = \frac{\text{percent}}{100}$$

\begin{itemize}[leftmargin=*]
    \item Put the \textbf{part} over the \textbf{whole} on the left side.
    \item Put the \textbf{percent} over $100$ on the right side.
    \item \textbf{Cross-multiply} and divide to solve for the unknown.
\end{itemize}

\medskip
\textbf{Example:} $18$ is $40\%$ of what number? $\to$ $\dfrac{18}{w} = \dfrac{40}{100}$ $\to$ $18 \times 100 = 40w$ $\to$ $w = 45$.
\end{quickReview}

% ── Warm-Up ──────────────────────────────────────────────────────────────────
\begin{practiceBox}[funGreen]{\faSun~Warm-Up}

\practiceHeader[funGreen]{Set Up the Proportion}
Write the proportion you would use to solve each problem. \textbf{Do not solve yet.}

\prob What is $35\%$ of $80$? \answerBlank[4cm]
\answerExplain{$\dfrac{n}{80} = \dfrac{35}{100}$}{The part is unknown ($n$), the whole is $80$, and the percent is $35$. Place them: $\frac{n}{80} = \frac{35}{100}$.}

\prob $24$ is what percent of $60$? \answerBlank[4cm]
\answerExplain{$\dfrac{24}{60} = \dfrac{p}{100}$}{The part is $24$, the whole is $60$, and the percent is unknown ($p$). Place them: $\frac{24}{60} = \frac{p}{100}$.}

\prob $45$ is $75\%$ of what number? \answerBlank[4cm]
\answerExplain{$\dfrac{45}{w} = \dfrac{75}{100}$}{The part is $45$, the whole is unknown ($w$), and the percent is $75$. Place them: $\frac{45}{w} = \frac{75}{100}$.}

\prob What is $90\%$ of $200$? \answerBlank[4cm]
\answerExplain{$\dfrac{n}{200} = \dfrac{90}{100}$}{The part is unknown ($n$), the whole is $200$, and the percent is $90$. Place them: $\frac{n}{200} = \frac{90}{100}$.}

\prob $14$ is what percent of $56$? \answerBlank[4cm]
\answerExplain{$\dfrac{14}{56} = \dfrac{p}{100}$}{The part is $14$, the whole is $56$, and the percent is unknown ($p$). Place them: $\frac{14}{56} = \frac{p}{100}$.}

\end{practiceBox}

% ── Practice A: Find the Part Using Proportions ──────────────────────────────
\begin{practiceBox}{Solving for the Part}

\practiceHeader{Set up a proportion and cross-multiply to find the part.}

\begin{multicols}{2}
\prob What is $45\%$ of $80$? \answerBlank[3cm]
\answerExplain{$36$}{Set up: $\frac{n}{80} = \frac{45}{100}$. Cross-multiply: $100n = 3{,}600$. Divide: $n = 36$.}

\prob What is $30\%$ of $150$? \answerBlank[3cm]
\answerExplain{$45$}{Set up: $\frac{n}{150} = \frac{30}{100}$. Cross-multiply: $100n = 4{,}500$. Divide: $n = 45$.}

\prob What is $12\%$ of $350$? \answerBlank[3cm]
\answerExplain{$42$}{Set up: $\frac{n}{350} = \frac{12}{100}$. Cross-multiply: $100n = 4{,}200$. Divide: $n = 42$.}

\prob What is $85\%$ of $60$? \answerBlank[3cm]
\answerExplain{$51$}{Set up: $\frac{n}{60} = \frac{85}{100}$. Cross-multiply: $100n = 5{,}100$. Divide: $n = 51$.}
\end{multicols}
\end{practiceBox}

% ── Practice B: Find the Percent Using Proportions ───────────────────────────
\begin{practiceBox}[funTeal]{Solving for the Percent}

\practiceHeader[funTeal]{Set up a proportion and cross-multiply to find the percent.}

\begin{multicols}{2}
\prob $21$ is what percent of $84$? \answerBlank[3cm]
\answerExplain{$25\%$}{Set up: $\frac{21}{84} = \frac{p}{100}$. Cross-multiply: $84p = 2{,}100$. Divide: $p = 25$, so $25\%$.}

\prob $36$ is what percent of $120$? \answerBlank[3cm]
\answerExplain{$30\%$}{Set up: $\frac{36}{120} = \frac{p}{100}$. Cross-multiply: $120p = 3{,}600$. Divide: $p = 30$, so $30\%$.}

\prob $54$ is what percent of $72$? \answerBlank[3cm]
\answerExplain{$75\%$}{Set up: $\frac{54}{72} = \frac{p}{100}$. Cross-multiply: $72p = 5{,}400$. Divide: $p = 75$, so $75\%$.}

\prob $8$ is what percent of $32$? \answerBlank[3cm]
\answerExplain{$25\%$}{Set up: $\frac{8}{32} = \frac{p}{100}$. Cross-multiply: $32p = 800$. Divide: $p = 25$, so $25\%$.}
\end{multicols}
\end{practiceBox}

% ── Practice C: Find the Whole Using Proportions ─────────────────────────────
\begin{practiceBox}[funPurple]{Solving for the Whole}

\practiceHeader[funPurple]{Set up a proportion and cross-multiply to find the whole.}

\begin{multicols}{2}
\prob $30$ is $40\%$ of what number? \answerBlank[3cm]
\answerExplain{$75$}{Set up: $\frac{30}{w} = \frac{40}{100}$. Cross-multiply: $40w = 3{,}000$. Divide: $w = 75$.}

\prob $63$ is $90\%$ of what number? \answerBlank[3cm]
\answerExplain{$70$}{Set up: $\frac{63}{w} = \frac{90}{100}$. Cross-multiply: $90w = 6{,}300$. Divide: $w = 70$.}

\prob $16$ is $20\%$ of what number? \answerBlank[3cm]
\answerExplain{$80$}{Set up: $\frac{16}{w} = \frac{20}{100}$. Cross-multiply: $20w = 1{,}600$. Divide: $w = 80$.}

\prob $48$ is $75\%$ of what number? \answerBlank[3cm]
\answerExplain{$64$}{Set up: $\frac{48}{w} = \frac{75}{100}$. Cross-multiply: $75w = 4{,}800$. Divide: $w = 64$.}
\end{multicols}
\end{practiceBox}

% ── Practice D: Word Problems ────────────────────────────────────────────────
\begin{practiceBox}[funBlue]{\faGlobe~Word Problems}

\wordProblem{In a bag of $200$ marbles, $35$ are blue. What percent of the marbles are blue?}{}
\answerExplain{$17.5\%$}{Set up the proportion: $\frac{35}{200} = \frac{p}{100}$. Cross-multiply: $200p = 3{,}500$. Divide: $p = 17.5$, so $17.5\%$ are blue.}

\wordProblem{A school has $520$ students. If $65\%$ ride the bus, how many students ride the bus?}{students}
\answerExplain{$338$ students}{Set up: $\frac{n}{520} = \frac{65}{100}$. Cross-multiply: $100n = 33{,}800$. Divide: $n = 338$ students ride the bus.}

\wordProblem{An athlete made $27$ free throws. This is $75\%$ of her attempts. How many did she attempt?}{attempts}
\answerExplain{$36$ attempts}{Set up: $\frac{27}{w} = \frac{75}{100}$. Cross-multiply: $75w = 2{,}700$. Divide: $w = 36$ total attempts.}

\end{practiceBox}

% ── Challenge ────────────────────────────────────────────────────────────────
\begin{challengeBox}
\prob A town has $3{,}200$ residents. Of them, $45\%$ are adults. Of the adults, $20\%$ volunteer at the library. Use proportions to find how many adults volunteer. \answerBlank[3cm]
\answerExplain{$288$ adults}{First find the number of adults: $\frac{n}{3{,}200} = \frac{45}{100}$, so $n = 1{,}440$. Then find volunteers: $\frac{v}{1{,}440} = \frac{20}{100}$, so $v = 288$ adults volunteer.}

\prob In a factory, $84$ items passed inspection. This was $70\%$ of the batch. Of the items that failed, $50\%$ can be repaired. How many items can be repaired? \answerBlank[3cm]
\answerExplain{$18$ items}{Find the total batch: $\frac{84}{w} = \frac{70}{100}$, so $w = 120$ items. Failed items: $120 - 84 = 36$. Can be repaired: $\frac{r}{36} = \frac{50}{100}$, so $r = 18$ items.}
\end{challengeBox}

\encouragement{Excellent! You now know how to solve any percent problem using proportions!}
